\documentclass[10pt,a4paper]{amsart}
\usepackage[margin=19mm]{geometry}
\usepackage{amsmath,amssymb,mathtools}
\usepackage[scr=rsfs,cal=euler]{mathalfa}
\usepackage{xfrac}
\usepackage[unicode,hidelinks,bookmarks=false]{hyperref}
\newcommand{\ie}{i.e.\ }
\newcommand{\cf}{cf.\ }
\newcommand{\resp}{respectively\ }
\newcommand{\Subsection}{\subsection}
\providecommand{\nobreakdash}{\nobreak\hbox{-}}
\setlength{\emergencystretch}{2em}
\multlinegap0pt
\usepackage{bm}
\usepackage{bbm}
\usepackage{dsfont}
\usepackage{xspace}
\usepackage{cancel}
\usepackage{mathtools}
\usepackage{amsthm}
\makeatletter
\@ifundefined{definition}{\newtheorem{definition}{Definition}}{}
\@ifundefined{lemma}{\newtheorem{lemma}[definition]{Lemma}}{}
\makeatother
\newcommand{\bbP}{\mathbb P}
\title[Bisection width, max-cut and internal partitions of 5-regula]{Bisection width, max-cut and internal partitions of 5-regular graphs}
\author{Endre Cs\'oka, Panna T\'imea Fekete, Zolt\'an L\'or\'ant Nagy, Levente Szemer\'edi}
\date{Source: arXiv:2509.08531v1 (CC BY 4.0)}
\begin{document}
\maketitle

\noindent\emph{Source and licence.} Bisection width, max-cut and internal partitions of 5-regular graphs. Version arXiv:2509.08531v1 (CC BY 4.0), distributed under the Creative Commons Attribution 4.0 International licence, as recorded in the arXiv licence metadata for this exact version and shown on the abstract page https://arxiv.org/abs/2509.08531v1. Full rights notice: licenses/arxiv-2509.08531-CC-BY-4.0.txt.\par\smallskip

\noindent\textit{Editorial scope.} Counted unit: the Lemma whose source label is \texttt{lemma:condindep} in the article (source0.tex lines 249--252, source label \texttt{lemma:condindep}), delivered together with the article's own complete original proof of it (source0.tex lines 254--271, one \texttt{proof} environment of 18 lines). No mathematics is written, repaired or re-derived: statement and proof are the article's own text line for line. Required context retained uncounted: none. Every definition, display and label of those lines is retained unchanged. Omitted from this package and not redistributed: 1--248, 253--253, 272--485. Nothing inside a retained excerpt is omitted or altered: every retained body is delivered exactly as the article's lines are written, so no omission receipt of any kind accompanies this package. Citations: the retained excerpts carry no citation command at all, so no bibliography entry of the article is transcribed and no reference is redistributed. Reference adaptation: none was needed and none was made; every reference inside the delivered unit resolves inside it, so no unresolved reference or citation is produced. Printed slips kept literal: no printed word, symbol, hypothesis or number of the counted statement or of its proof was altered. Layout and class: the article's own class is \texttt{article} with packages including babel, geometry, amsmath, amsthm, amssymb, mathtools, graphicx, hyperref, makecell, hhline, comment, listings, bbold, soul; the wrapper is the lane's pinned \texttt{amsart} editorial wrapper at 10pt on A4, which restates the theorem-like environments the retained text needs and adds the article's own macro definitions that the retained text needs (\texttt{\textbackslash bbP}), with extra wrapper packages (bm, bbm, dsfont, xspace, cancel, mathtools). The delivered numbering is the unit's own. Assets: no retained span contains a figure environment, a graphics-inclusion command, a PDF-page inclusion, a table or a TikZ picture; no image file is copied into this package, the package declares no asset dependency and no image file is redistributed with it. Licence: version arXiv:2509.08531v1 of this article is distributed under the Creative Commons Attribution 4.0 International licence, as recorded in the arXiv licence metadata for this exact version and shown on the abstract page https://arxiv.org/abs/2509.08531v1. A full rights notice is delivered at licenses/arxiv-2509.08531-CC-BY-4.0.txt. Characters: every retained excerpt is pure ASCII, so the delivered engine, XeLaTeX with its default Latin Modern text font, reports no missing character.
\begin{lemma}
\label{lemma:condindep}
 Let $w_1$ and $w_2$ be two adjacent vertices of the $d$-regular infinite tree. If we remove the edge $w_1w_2$ the graph falls into two connected components $S_1$ and $S_2$. Let $A_1\subset S_1$ and $A_2\subset S_2$ be any finite set of vertices. Then for any step $t$, the colorings $C_t(A_1)$ and $C_t(A_2)$ are independent conditioned on the event $(C_{t-1}(w_1)="U")\cap(C_{t-1}(w_2)="U"),$ meaning that if there is an edge with both of its endpoint uncolored then the two sides of the edge are independent.
\end{lemma}

\begin{proof}
 With $i\in\{1,2\}$ and $0\leq s\leq t$ integer, define the set $A_i^s$ recursively as follows: $A_i^t=A_i$ and $A_i^s=N[A_i^{s+1}]\setminus\{w_{2-i}\}$ if $s<t$, i.e. the closed neighborhood of $A^{s+1}_i$ apart from the vertex $w_{2-i}$.

We prove the following stronger statement: the colorings $C_s(A_1^s)$ and $C_s(A_2^s)$ are independent conditioned on the event $(C_{t-1}(w_1)="U")\cap(C_{t-1}(w_2)="U").$
    
 It is true for $s=0$: in the beginning every vertex is uncolored which means that $C_0$ is constant. So $C_0(A_1^0)$ and $C_0(A_2^0)$ are also constant random variables, thus they are independent even when conditioned on any event.
    
 Now suppose that the statement is true for some $s<t$. We want the conditional independence of $C_{s+1}(A_1^{s+1})$ and $C_{s+1}(A_2^{s+1})$. With $i\in\{1,2\}$, let $T_i^{s+1}:A_i^{s+1}\to\{"U","R","B"\}$ be an arbitrary coloring on $A_i^{s+1}$. During the following calculation we use the following short hand notation for the event we condition on: $B_t=(C_{t-1}(w_1)="U")\cap(C_{t-1}(w_2)="U").$ Note that $B_t$ also implies $B_1,\dots,B_{t-1}.$ Now we can prove the conditional independence.
\begin{equation*}
\begin{split}
 \bbP\left(\bigcap_{i\in\{1,2\}}\{C_{s+1}(A_i^{s+1})=T_i^{s+1}\}\middle|B_t\right)&=\sum_{\substack{R^s_1\in\mathcal T^c_{A_1^s},\\R^s_2\in\mathcal T^c_{A_2^s}}}\bbP\left(\bigcap_{i\in\{1,2\}}\{C_{s+1}(A_i^{s+1})=T_i^{s+1},C_s(A_i^s)=R_i^s\}\middle| B_t\right)\\
 &=\sum_{\substack{R^s_1\in\mathcal T^c_{A_1^s},\\R^s_2\in\mathcal T^c_{A_2^s}}}\prod_{i\in\{1,2\}} \bbP\left(C_{s+1}(A_i^{s+1})=T_i^{s+1},C_s(A_i^s)=R_i^s\middle| B_t\right)\\
 &=\prod_{i\in\{1,2\}} \bbP\left(C_{s+1}(A_i^{s+1})=T_i^{s+1}\middle| B_t\right).\\
\end{split}
\end{equation*}

The first and the last equality is true because the sums run through disjoint events. For the equation in the middle, we need to understand what happens to a fixed vertex at one step. Based on the colors of its neighbors it either stays the same or using some randomness it changes color. The random choice of the vertices are independent of each other. So the behavior of one vertex at one step only depends on the state of its one-neighborhood and some randomness of its own. By fixing the coloring after the step $s$ as $C_s(A_i^s)=R_i^s$, we only need to check if the one-neighbors overlap or not. Apart from $w_1$ and $w_2$ any vertex $v_i\in A_i$ has $N[v_i]\subset S_i$. But $C_{s}(w_1)=C_{s}(w_2)="U"$ is forced by the condition $B_t$. So the colorings $C_{s+1}(A_1^{s+1})$ and $C_{s+1}(A_2^{s+1})$ are independent under the condition $B_t$ when we know the colorings $C_s(A_i^s)$.
\end{proof}

\end{document}
